% Generated by `taskgen course examples` from factoring.lenstra (seed 3).
% Do not edit: change examples.yaml and regenerate.
\begin{example}[алгоритм Ленстры, неудача с $\gcd = N$]\label{ex:ecm-lenstra-gcd-n}
Найти нетривиальный делитель числа $N = 221$ алгоритмом Ленстры. Кривая $E_1\colon y^2 = x^3 + 97x + b_1$ над $\mathbb{Z}/N\mathbb{Z}$, где $b_1$ выбрано так, что точка $P_1 = (60, 144)$ лежит на $E_1$. Вычислять $Q_k = k\,Q_{k-1}$ ($Q_1 = P_1$, т.\,е. $Q_k = k!\,P_1$), $k = 2, 3, \dots$, удвоением-сложением по модулю $N$, пока не встретится необратимый знаменатель. Если алгоритм на $E_1$ не даст нетривиального делителя, взять кривую $E_2\colon y^2 = x^3 + 145x + b_2$ с точкой $P_2 = (205, 5)$.

\par\noindent\textit{Решение.} Формулы над $\mathbb{Z}/N\mathbb{Z}$: $\lambda = (3x_1^2 + a)/(2y_1)$ (удвоение), $\lambda = (y_2 - y_1)/(x_2 - x_1)$ (сложение), $x_3 = \lambda^2 - x_1 - x_2$, $y_3 = \lambda(x_1 - x_3) - y_1$. Деление на $d$ --- умножение на $d^{-1} \bmod N$ (расширенный алгоритм Евклида для $(N, d)$); если $\gcd(d, N) > 1$, обратного нет, и $\gcd(d, N)$ --- делитель $N$. Умножение на $k$ --- слева направо по двоичной записи $k$: $R \gets Q$, для каждого следующего бита $R \gets 2R$, при бите $1$ ещё $R \gets R + Q$.\par\medskip\noindent\textbf{Кривая $E_1$.} $b = y_0^2 - x_0^3 - a x_0 = 144^2 - 60^3 - 97 \cdot 60 \equiv 26\pmod{221}$, кривая $E_1\colon y^2 = x^3 + 97x + 26$; $4a^3 + 27b^2 \equiv 123$, $\gcd(123, 221) = 1$ (кривая невырождена по модулю обоих простых делителей). $Q_1 = P_1 = (60, 144)$; строки $k$ вычисляют $Q_k = k\,Q_{k-1}$, $Q = Q_{k-1}$, $m$ --- кратность точки $R = mP$ после шага, $d$ --- знаменатель:\[{\small\begin{array}{c|l|r|r|r|r|c}k & \text{шаг} & m & d & d^{-1} & \lambda & R = mP_1 \\ \hline 2 = 10_2 & 2R & 2 & 67 & 33 & 34 & (152, 43) \\ \hline 3 = 11_2 & 2R & 4 & 86 & 18 & 49 & (108, 124) \\  & R+Q & 6 & 44 & 216 & 184 & (4, 6) \\ \hline 4 = 100_2 & 2R & 12 & 12 & 129 & 141 & (204, 82) \\  & 2R & 24 & 164 & 31 & 49 & (4, 215) \\ \hline 5 = 101_2 & 2R & 48 & 209 & 92 & 80 & (204, 139) \\  & 2R & 96 & 57 & 190 & 172 & (4, 6) \\  & R+Q & 120 & 0 & \gcd = 221 & \text{---} & \text{---} \\ \end{array}}\]\par Знаменатель $d = 0$ при $k = 5$ необратим: $\gcd(0, 221) = 221$ (остатки алгоритма Евклида: $221, 0$). Объяснение: $\#E(\mathbb{F}_{13}) = 10$, $\operatorname{ord}(P_1 \bmod 13) = 2 \cdot 5$; $\#E(\mathbb{F}_{17}) = 15$, $\operatorname{ord}(P_1 \bmod 17) = 3 \cdot 5$; на последнем шаге вычисляется точка $120P_1$ ($120 = 2^{3} \cdot 3 \cdot 5$), и она равна $\mathcal{O}$ по модулю обоих простых сразу.\par Делитель тривиален ($\gcd = N$): порядки $P_1$ по модулю обоих простых делят $120$. Неудача --- берём вторую кривую.\par\medskip\noindent\textbf{Кривая $E_2$.} $b = y_0^2 - x_0^3 - a x_0 = 5^2 - 205^3 - 145 \cdot 205 \equiv 32\pmod{221}$, кривая $E_2\colon y^2 = x^3 + 145x + 32$; $4a^3 + 27b^2 \equiv 185$, $\gcd(185, 221) = 1$ (кривая невырождена по модулю обоих простых делителей). $Q_1 = P_2 = (205, 5)$; строки $k$ вычисляют $Q_k = k\,Q_{k-1}$, $Q = Q_{k-1}$, $m$ --- кратность точки $R = mP$ после шага, $d$ --- знаменатель:\[{\small\begin{array}{c|l|r|r|r|r|c}k & \text{шаг} & m & d & d^{-1} & \lambda & R = mP_2 \\ \hline 2 = 10_2 & 2R & 2 & 10 & 199 & 25 & (215, 187) \\ \hline 3 = 11_2 & 2R & 4 & 153 & \gcd = 17 & \text{---} & \text{---} \\ \end{array}}\]\par Знаменатель $d = 153$ при $k = 3$ необратим: $\gcd(153, 221) = 17$ (остатки алгоритма Евклида: $221, 153, 68, 17, 0$). Объяснение: $\#E(\mathbb{F}_{13}) = 16$, $\operatorname{ord}(P_2 \bmod 13) = 2^{3}$; $\#E(\mathbb{F}_{17}) = 24$, $\operatorname{ord}(P_2 \bmod 17) = 2^{2}$; на последнем шаге вычисляется точка $4P_2$ ($4 = 2^{2}$), и она равна $\mathcal{O}$ по модулю $17$, но не по модулю $13$.\par Делитель $17$, $N / 17 = 13$.\[N = 13 \cdot 17 = 221.\]
\end{example}
